\documentclass[11pt,a4paper]{article}

\usepackage[utf8]{inputenc}
\usepackage[T1]{fontenc}
\usepackage[english]{babel}
\usepackage{amsmath,amssymb,amsthm}
\usepackage{graphicx}
\usepackage{geometry}
\usepackage{hyperref}
\usepackage{xcolor}
\usepackage{url}
\usepackage{booktabs}
\usepackage{authblk}
\usepackage{caption}
\usepackage{enumitem}

\geometry{margin=2.5cm}
\hypersetup{
    colorlinks=true,
    linkcolor=blue!50!black,
    citecolor=blue!50!black,
    urlcolor=blue!50!black,
}

\title{\textsc{Dynamoscope}: Treating Video as a Trajectory in State Space \\
       \large A Unified Framework for Multi-Scale, Multi-Modal Emergence Analysis}

\author{Zaccarie Aglas}
\affil{Independent Researcher, Paris, France \\ \texttt{zaccarie@zaccarie.fr}}

\date{\today}

\begin{document}

\maketitle

\begin{abstract}
We introduce \textsc{Dynamoscope}, an experimental research tool that
re-frames video analysis as the study of a single trajectory in a
high-dimensional state space. Rather than treating frames as independent
images, we encode each frame with a pretrained visual model and analyze
the resulting sequence as a dynamical system. The pipeline integrates
sixteen analytical modules ranging from classical dynamical systems
metrics (Lyapunov exponent, recurrence quantification, correlation
dimension) to topological data analysis (persistent homology),
information theoretic causality (Granger, transfer entropy,
convergent cross-mapping), sparse equation discovery (SINDy), slow
feature extraction (SFA), latent prediction with counterfactual
rollouts (predictive coding / JEPA-style imagination), and
unsupervised cluster discovery (HDBSCAN) with optional CLIP-based
semantic labeling. Five encoders (ResNet50, ViT-B/16, DINOv2-S, CLIP
ViT-B/32, and VideoMAE-base) and an aligned audio modality
(wav2vec2) allow direct comparison of how different pretraining
objectives---including a natively temporal one---shape the
recovered trajectory geometry. We argue that the gaps and clusters seen in
non-linear projections such as UMAP are not noise but structural
detectors operating at a chosen scale, and we propose an interactive
parameter optimization layer (random and Bayesian) that lets the user
discover representations under explicit faithfulness, continuity, and
spread objectives. Five view modes (phase-space, time-helix, Takens,
polar, audio, SFA) and five color modes provide epistemic
triangulation. The tool runs on consumer hardware (Apple MPS / NVIDIA
CUDA / CPU), supports live webcam capture, persists experiments in
a SQLite registry, and exports standalone HTML reports. We validate
the pipeline on synthetic systems (Lorenz, R\"ossler, H\'enon, double
pendulum) where ground-truth equations are recoverable to $R^2 > 0.99$
via SINDy, and demonstrate its qualitative behavior on real video.
\textsc{Dynamoscope} is intended as a microscope for emergent
structure in time-varying visual signals.
\end{abstract}

\textbf{Keywords:} video understanding; dynamical systems; persistent
homology; UMAP; self-supervised learning; causal emergence; world
models; Takens embedding; Koopman operator; SINDy; slow feature
analysis.

\section{Introduction}

A video is, simultaneously, a sequence of images and a sample
trajectory of an underlying process. The first reading dominates
contemporary computer vision: deep networks classify, detect, segment,
or caption frames, typically treating temporal context as an
afterthought to be handled by recurrent or attention layers. The
second reading, native to physics and applied dynamics, is comparatively
underdeveloped despite the fact that virtually every interesting video
encodes a process: a body in motion, a scene in transition, a system
evolving toward or away from attractors.

\textsc{Dynamoscope} is built around the second reading. We treat a
video $\mathcal{V}$ as a discrete-time sample of a continuous
process $\Phi: \mathbb{R}_+ \to \mathcal{X}$, where $\mathcal{X}$
is some latent space of \emph{perceptual states}. By choosing a
visual encoder $f_\theta: \mathbb{R}^{H\times W\times 3} \to
\mathbb{R}^D$, we obtain a sequence of state vectors
$z_t = f_\theta(x_t)$, $t=1,\dots,N$, which can then be subjected to
the full battery of techniques developed for studying dynamical
systems: phase-space reconstruction, attractor analysis,
Lyapunov exponents, topology of trajectories, predictive modeling,
and so on.

The framework is intentionally encoder-agnostic. Different encoders
produce different state spaces, and one of our main empirical points
is that the geometry of the recovered trajectory depends
\emph{strongly} on $f_\theta$. We argue this is a feature, not a bug:
the encoder defines what is perceptually salient, and comparing
encoders reveals their inductive biases concretely.

Our contributions are:
\begin{enumerate}[leftmargin=*,topsep=2pt,itemsep=2pt]
    \item A unified pipeline that ingests a video (or a synthetic
    dynamical system, or a live webcam stream), produces a
    high-dimensional latent trajectory, projects it to a visualizable
    3D embedding, and feeds it to a suite of sixteen analytical
    modules.
    \item An encoder zoo (ResNet50~\cite{he2016deep},
    ViT-B/16~\cite{dosovitskiy2020vit},
    DINOv2~\cite{oquab2023dinov2},
    CLIP~\cite{radford2021clip}) plus an aligned audio modality
    (wav2vec2~\cite{baevski2020wav2vec2}), enabling direct
    comparison of representational biases.
    \item A parameter-tuning layer for the visual projection step
    (UMAP~\cite{mcinnes2018umap}, PCA, Isomap~\cite{tenenbaum2000isomap})
    using random search and Bayesian
    optimization~\cite{snoek2012bayesian} under an explicit composite
    objective (faithfulness, continuity, spread).
    \item Five view modes (phase-space, time-helix, Takens delay
    embedding, polar, audio trajectory, slow-feature trajectory) and
    five color modes (time, two velocity variants, cluster, prediction
    surprise) supporting epistemic triangulation across representations.
    \item Predictive modeling in latent space via a residual MLP, with
    deterministic and counterfactual rollouts that visualize a
    ``cone of imagined futures'', directly analogous to active
    inference~\cite{friston2010free} and JEPA-style world
    models~\cite{lecun2022path}.
    \item A composite ``video DNA'' score that aggregates seven
    dimensions (chaos, topology, complexity, spectral, structure,
    causality, predictability) into a single interpretable index
    with classification.
    \item Persistent storage, multi-video comparison via shared
    UMAP fit, standalone HTML reports, and a fully interactive
    Three.js front-end with hover tooltips, raycaster picking, and
    live parameter sliders.
\end{enumerate}

The pipeline runs in real time on consumer hardware (Apple Silicon
M-series via MPS, NVIDIA CUDA, or CPU). All code is in Python (FastAPI
backend) and a single-file vanilla JS front-end.

\section{Related Work}

\paragraph{Dynamical systems on observed signals.}
The treatment of a time series as a partial observation of a
dynamical system is classical. Takens' embedding
theorem~\cite{takens1981} establishes that delay-coordinate
reconstruction $(s(t), s(t+\tau), \dots, s(t+(m-1)\tau))$ is
diffeomorphic to the true attractor of the underlying flow under
generic conditions. Rosenstein, Collins, and De
Luca~\cite{rosenstein1993} provided a practical estimator of the
largest Lyapunov exponent, while Grassberger and
Procaccia~\cite{grassberger1983} defined the correlation dimension
that we use as a proxy of fractal complexity. Recurrence
quantification analysis~\cite{marwan2007recurrence} extracts
diagonal-line statistics from recurrence plots. We apply all four
to the latent trajectory.

\paragraph{Topological data analysis.}
Persistent homology~\cite{edelsbrunner2002topological,
carlsson2009topology} quantifies the multi-scale topological
features (connected components, loops, voids) of a point cloud or
sublevel-set filtration. Recent work on time-series topology
includes sliding-window persistence and persistence
landscapes~\cite{perea2015sw1pers}. We use the \texttt{ripser}
implementation~\cite{tralie2018ripser} on both the global trajectory
and sliding windows.

\paragraph{Self-supervised visual representations.}
The shift from supervised ImageNet pretraining
(ResNet50~\cite{he2016deep}, ViT-B~\cite{dosovitskiy2020vit}) to
self-supervised contrastive (CLIP~\cite{radford2021clip}) and
self-distillation (DINO, DINOv2~\cite{oquab2023dinov2}) objectives
has produced models with substantially different feature geometry:
self-supervised features tend to vary more smoothly with continuous
visual transformations~\cite{caron2021dino}, which we observe
empirically in the recovered trajectory smoothness.

\paragraph{World models and predictive coding.}
The view that biological and artificial agents maintain an internal
generative model of their sensory stream
\cite{friston2010free, rao1999predictive} motivates our predictive
module. JEPA~\cite{lecun2022path} explicitly advocates
prediction in latent rather than pixel space, which is exactly
what our residual MLP does. Dreamer~\cite{hafner2020dream} and
related recurrent state-space models maintain a similar latent
forward model. Our counterfactual rollout is closely related to
imagination-based planning in active inference~\cite{kaplan2018planning}.

\paragraph{Sparse equation discovery and slow features.}
SINDy~\cite{brunton2016sindy} recovers governing equations of a
dynamical system from time series via sparse regression on a
polynomial library. Slow Feature Analysis~\cite{wiskott2002slow}
extracts the linear combinations of inputs that vary most slowly
in time, and has been related to the principle of slowness in
biological vision~\cite{berkes2005slow}. Both are integrated
in our pipeline as alternative representational lenses.

\paragraph{Information-theoretic causality and emergence.}
Granger causality~\cite{granger1969} tests whether the past of one
series helps predict another. Transfer entropy~\cite{schreiber2000te}
generalizes this to a non-parametric, information-theoretic measure.
Convergent cross-mapping~\cite{sugihara2012ccm} uses Takens-style
manifold reconstruction to detect causality in deterministic
non-linear systems. Causal emergence~\cite{hoel2013quantifying}
quantifies when a coarse-grained description carries more
effective information than the micro-scale, formalizing a notion of
``emergent'' macro-causation that we approximate.

\paragraph{UMAP and projection caveats.}
UMAP~\cite{mcinnes2018umap} preserves local topological structure
but not global distances. Gaps and clusters in a UMAP projection are
therefore artifacts of the chosen $k$-neighborhood scale, not
distance-preserving features of the original space. We expose this
explicitly through tunable parameters and an objective-based
auto-tuning layer.

\section{Methods}

\subsection{Pipeline overview}

Given a video file or a webcam stream, frames are sampled uniformly
(by default 150 frames) and resized to $224 \times 224$. Each frame
is passed through the selected encoder, producing a sequence
$\{z_t\}_{t=1}^N \subset \mathbb{R}^D$. If the source carries
audio, the soundtrack is also extracted and aligned: for each
sampled frame index $t$, a one-second window of audio centered on
the corresponding timestamp is encoded by wav2vec2 and pooled to
yield $a_t \in \mathbb{R}^{768}$. The visual sequence is then
projected to $\mathbb{R}^3$ via UMAP (default), PCA, or Isomap, and
optionally smoothed by a temporal moving average. The 3D
coordinates are normalized to the cube $[-1,1]^3$ for stable
rendering. All downstream analyses are computed on either the 3D
projection, the original $D$-dimensional latents, or both.

\subsection{Encoders}

We provide four pretrained encoders, summarized in
Table~\ref{tab:encoders}. Their outputs differ in dimensionality,
geometry, and inductive bias.

\begin{table}[ht]
\centering
\begin{tabular}{lrrl}
\toprule
Encoder & Dim & Params & Pretraining objective \\
\midrule
ResNet50 V2~\cite{he2016deep}    & 2048 & 25\,M  & Supervised ImageNet \\
ViT-B/16~\cite{dosovitskiy2020vit} & 768  & 86\,M  & Supervised ImageNet \\
DINOv2-S/14~\cite{oquab2023dinov2}& 384  & 22\,M  & Self-supervised distillation \\
CLIP ViT-B/32~\cite{radford2021clip} & 512  & 151\,M & Contrastive image--text \\
VideoMAE-base~\cite{tong2022videomae} & 768  & 86\,M  & Self-supervised, 16-frame clips (Kinetics-400) \\
\bottomrule
\end{tabular}
\caption{Encoders shipped with \textsc{Dynamoscope}. VideoMAE is the only
encoder that natively integrates temporal context: each output feature
aggregates a 16-frame window centered on the target frame.}
\label{tab:encoders}
\end{table}

For audio: \texttt{facebook/wav2vec2-base}~\cite{baevski2020wav2vec2}
(94\,M parameters, 768-dim output, 1-second context windows centered
on each frame timestamp).

\subsection{Dimensional reduction}

We default to UMAP with cosine metric, $n_{\text{neighbors}}=15$,
$\min_{\text{dist}}=0.1$. PCA is offered as a linear, distance-preserving
alternative, and Isomap as a $k$-NN geodesic compromise. We
emphasize that any non-linear projection trades distance fidelity
for topological clarity. For this reason, all downstream metrics
that depend on metric structure (Lyapunov, correlation dimension,
recurrence) can be optionally computed on either the projected 3D
coordinates or the original $D$-dimensional latents.

\subsection{Dynamical analysis}

\paragraph{Lyapunov exponent.} We use the Rosenstein
estimator~\cite{rosenstein1993}: for each point in the trajectory,
the nearest non-temporal neighbor is identified, and the divergence
$d(k)$ of their forward orbits is averaged over the dataset. The
slope of $\log d(k)$ versus $k$ over a fitting window yields the
largest Lyapunov exponent.

\paragraph{Correlation dimension.} Estimated by Grassberger--Procaccia
\cite{grassberger1983}: the cumulative pairwise distance distribution
$C(\varepsilon) \propto \varepsilon^d$ is fitted by linear regression
in log--log coordinates.

\paragraph{Recurrence quantification.} We threshold the pairwise
distance matrix at the 10th percentile to obtain a binary recurrence
plot. Three statistics summarize it: \emph{RR} (recurrence rate),
\emph{DET} (determinism, fraction of recurrent points lying on
diagonal lines of length $\geq 2$), and \emph{LAM} (laminarity, the
analogous vertical-line statistic).

\paragraph{Takens delay embedding.} We expose a view mode that, given
the first principal coordinate as a scalar observable $s(t)$,
reconstructs a 3D phase portrait $(s(t), s(t+\tau), s(t+2\tau))$
with user-tunable delay $\tau$. For Lorenz, this recovers a
qualitatively faithful butterfly from a single scalar.

\subsection{Spectral analysis and Koopman}

We compute the exact DMD~\cite{tu2014dmd} on the latent trajectory.
The linear operator approximating $z_{t+1} \approx A z_t$ is
constructed via truncated SVD; its eigendecomposition yields modes
whose eigenvalues lie on or near the unit circle for stable
oscillations and decay inside it for transients. The continuous
frequencies are $\omega = \log(\lambda)/\Delta t$. The unit-circle
visualization makes the Koopman approximation immediately legible.
We complement DMD with a multi-scale entropy ladder (Shannon entropy
of histogrammed latent dimensions across temporal coarse-grainings
$s \in \{1,2,4,8,16,32\}$) and a power-law spectral slope obtained
from log--log fitting of $|\text{FFT}(z)|^2$.

\paragraph{Wavelet decomposition.} For each latent dimension we
optionally compute a Discrete Wavelet Transform (DWT) via
PyWavelets, with a choice of mother wavelets (Daubechies-4 by
default, Haar, Symlet-8, Coiflet-5). The energy per scale
$E_j = \sum_k |cD_{j,k}|^2$ summarizes the temporal multi-scale
structure: a dominant fine scale indicates rapid oscillations,
while energy concentrated at coarse scales suggests slow modulation.
This is complementary to DMD (frequency-domain) and entropy ladder
(distribution-shape) signatures. A proper Mallat scattering
network~\cite{mallat2012scattering} would yield Lipschitz-stable
translation-invariance; we view the DWT decomposition as a
lightweight surrogate.

\subsection{Causality}

We provide three complementary tests pairwise on the dominant
latent components:
\begin{itemize}[leftmargin=*]
    \item \textbf{Granger}~\cite{granger1969}: compares the residual
    sum of squares of an autoregression with and without the past of
    a candidate cause.
    \item \textbf{Transfer entropy}~\cite{schreiber2000te}:
    information-theoretic measure
    $\mathrm{TE}_{y\to x} = I(x_t; y_{t-\tau} \mid x_{t-\tau})$.
    Robust to non-linear coupling.
    \item \textbf{Convergent cross-mapping (CCM)}~\cite{sugihara2012ccm}:
    tests whether the shadow manifold reconstructed from one
    variable contains information about another, indicating
    deterministic coupling.
\end{itemize}
Edge agreement across the three tests is taken as a robustness
indicator.

\subsection{Topology}

We compute persistent homology of dimensions 0 and 1 on the 3D
trajectory using \texttt{ripser}~\cite{tralie2018ripser}. The
resulting birth--death diagram is rendered, along with persistence
entropy~\cite{rucco2017persistence} as a scalar summary. A sliding-window
variant tracks the number and stability of $H_1$ cycles over time,
acting as a temporal topological signature.

\subsection{Sparse equation discovery}

SINDy~\cite{brunton2016sindy} is applied directly to the (denoised)
trajectory. We build a polynomial library up to degree 2 and use the
sequential thresholded least-squares algorithm to identify a sparse
coefficient matrix. For Lorenz with $(\sigma, \rho, \beta) =
(10, 28, 8/3)$, SINDy recovers the canonical equations with
coefficients within 5\% of ground truth (Section~\ref{sec:results}).

\subsection{Slow Feature Analysis}

SFA~\cite{wiskott2002slow} is implemented via the generalized
eigenvalue formulation. After centering and whitening, we
diagonalize the covariance of the temporal derivative; the smallest
eigenvalues correspond to the slowest features. These can be used
directly as a 3D coordinate system (\emph{SFA view mode}), providing
an alternative projection whose axes are interpretable as
``temporal invariants''.

\subsection{Emergence}

We approximate causal emergence~\cite{hoel2013quantifying} by
computing the effective information
$\mathrm{EI}(s) = I(z_t^{(s)}; z_{t+1}^{(s)})$ at multiple
coarse-graining scales $s$. The maximum gain
$\phi_{\text{id}} \approx \max_s \mathrm{EI}(s) - \mathrm{EI}(1)$
over micro-scale serves as a coarse indicator of when macro
descriptions are more predictive than the underlying micro
trajectory.

\subsection{Latent prediction and counterfactual rollouts}

We train a residual MLP $g_\phi$ such that
$z_{t+1} \approx z_t + g_\phi(z_t)$ via Adam minimization of MSE in
standardized latent space. Per-frame surprise is defined as the
prediction error, providing a color mode that highlights unexpected
events. The trained predictor enables iterative rollout: starting
from $z_{t_0}$, we generate $z_{t_0+1}, z_{t_0+2}, \dots$ by repeated
application. Step-wise error against the true trajectory yields an
empirical predictive horizon. Adding Gaussian noise to the starting
state and rolling out several samples produces a ``cone of imagined
futures'', visually conveying the local divergence rate (an
empirical analog of the largest Lyapunov exponent).

\subsection{Clustering and transitions}

HDBSCAN~\cite{campello2013hdbscan} is run on $\ell_2$-normalized
latents to discover modes without specifying $k$. The resulting
labels can be used as a color mode, and the temporal sequence of
labels yields a transition matrix, visualized as a directed graph
where node size encodes time spent in each cluster and edge
thickness encodes transition count. Optionally, each cluster is
labeled by CLIP-based comparison of its centroid feature against a
user-supplied vocabulary, yielding automatic semantic tags.

\subsection{Composite ``video DNA''}

We define a 7-axis scorecard (chaos, topology, complexity, spectral
slope, structure, causality, predictability), normalized to $[0,1]$,
and a weighted sum yielding a single composite score in $[0, 100]$.
A simple rule-based classifier assigns a textual label (e.g.,
\emph{chaotic with periodic structure}, \emph{structured / clustered
/ predictable}) based on threshold combinations. The score is
intended as a quick comparative index across runs, not as a
definitive measure.

\section{Results}
\label{sec:results}

\subsection{Validation on synthetic systems}

The synthetic system zoo (Lorenz, R\"ossler, Van der Pol, H\'enon,
logistic map, double pendulum) provides ground truth against which
the pipeline can be validated. On Lorenz with default parameters,
SINDy recovers the equations
\[
\dot x \approx -9.84 x + 9.97 y, \quad
\dot y \approx 25.87 x - 0.96 xz, \quad
\dot z \approx 0.996 xy - 2.26 z
\]
with $R^2 = 0.997$, compared to ground truth coefficients
$(-10, 10, 28, -1, -1, 1, -8/3)$. (The $-y$ term in $\dot y$ falls
below our 0.5 threshold and is suppressed; the constant in
$\dot z$ reflects integration drift.)

The Koopman spectrum on the same trajectory contains three modes
with $|\lambda| \approx 1$ and continuous frequencies near
$\pm 0.81$\,Hz, consistent with the dominant orbital period of the
butterfly. Multi-scale entropy is essentially flat across
coarse-grainings $s \in \{1,\dots,32\}$, indicating scale-invariant
chaotic structure. The spectral slope is approximately $-3$, in
line with the Kolmogorov regime expected for low-dimensional chaos.
Persistent homology recovers 5--7 stable $H_1$ cycles corresponding
to the closed orbits around each butterfly wing.

\subsection{Encoder effect on real video}

On a 6-second continuous Mandelbrot zoom (90 sampled frames), the
five available encoders produce qualitatively different latent
trajectories. ResNet50 yields a piecewise zigzag, reflecting the
discrete jumps activated as the zoom traverses the boundaries of
ImageNet category responses. ViT-B/16 is similar. DINOv2-S produces
a smooth monotonic curve, with mean cosine velocity
$\approx 0.009$ per frame --- roughly 2.6$\times$ smoother than
ResNet50. CLIP sits between the two. VideoMAE-base, the only
encoder that integrates a 16-frame temporal context per output,
produces a trajectory three orders of magnitude smoother than
DINOv2 (mean cosine velocity $\approx 5\times 10^{-5}$). This is
consistent with its training objective: by reconstructing masked
spatio-temporal patches across 16-frame windows, VideoMAE learns
features that are explicitly invariant to local motion, yielding an
ultra-smooth latent stream for any continuous transformation. The
trade-off is computational: per-frame encoding is approximately
$10\times$ slower than single-frame encoders, which limits live
applications. This direct comparison across five encoders is, in
our view, one of the most pedagogically useful features of the
tool.

\subsection{UMAP gaps as scale-dependent structure}

The UMAP projection of the same video exhibits visible gaps even
when the underlying latent perceptual velocity is uniform. We
provide three independent diagnostics:
\begin{itemize}[leftmargin=*,topsep=2pt]
    \item Comparing velocity in the projection vs.\ velocity in
    original $D$-space confirms that projection gaps are largely
    artifacts of UMAP's repulsive forces.
    \item Switching to PCA recovers a continuous trajectory whose
    distances proportionally match the original ones, at the cost of
    visual cluster separation.
    \item Tuning $n_{\text{neighbors}}$ from 5 to 50 sweeps the
    projection from a fragmented local microstructure to a globally
    integrated macro-structure, illustrating that the gaps are
    detectors of structure \emph{at a chosen scale} rather than
    intrinsic features.
\end{itemize}

\subsection{Auto-tuning the projection}

Twenty random-search trials on Mandelbrot under our composite
objective return $n_{\text{neighbors}}=19$, $\min_{\text{dist}}=0.27$,
metric=\texttt{correlation}, smoothing=1 with composite score
$0.927$. Bayesian optimization with the same budget converges to
slightly different parameters and similar score; the difference is
within run-to-run noise for small budgets, consistent with
\cite{bergstra2012random}. For budgets above $\sim$25 trials,
Bayesian optimization is preferable.

\subsection{Counterfactual cones and empirical Lyapunov}

Starting a counterfactual cone from frame 25 of Mandelbrot with
$H=20$ steps and 10 samples at $\sigma=0.08$, we observe a maximum
inter-sample 3D distance of $\sim 1.07$ at the final step --- about
half of the available range. The divergence rate is broadly
consistent with the step-wise error growth observed in a
deterministic rollout, which we interpret as an empirical
Lyapunov-like quantity in the latent space.

\subsection{Cross-modal causality}

On a 6-second video with a constant 440\,Hz sine tone added to a
Mandelbrot zoom, the cross-modal transfer entropy matrix between
visual and audio principal components shows weak but non-trivial
coupling. This is expected: the audio is constant in content but
not in our pooling, and the visual evolves continuously. For
real-world videos with music or speech aligned to events, this
analysis surfaces the synchronization patterns directly.

\section{Discussion}

\subsection{The encoder is the lens}

A recurring theme throughout the design is that the encoder is not a
neutral observer. ResNet50 pretrained on ImageNet sees the world
through a thousand discrete categories; it produces feature
activations that jump as a continuous visual stream crosses category
boundaries. DINOv2, trained without labels, produces a smoother,
more continuous geometry. CLIP, trained on weakly aligned
captions, lies between these two. We argue that this is not a defect
but a fact about visual perception by computational means:
\emph{there is no view from nowhere}. \textsc{Dynamoscope}'s value is
to make this concrete and comparable.

\subsection{Projection gaps as scale detectors}

A common misreading of UMAP visualizations is that gaps between
clusters indicate genuine ``empty'' regions in the original space.
Our experiments show that for continuous transformations like a
zoom, the latent path is in fact smooth in the original
$D$-dimensional space; UMAP introduces gaps because its loss
function explicitly pushes non-neighbors apart at a scale set by
$n_{\text{neighbors}}$. This is informative if read correctly: a
gap reports that the algorithm \emph{decided}, at this scale, that
two neighborhoods are topologically distinct. Reading it as
``there is empty space in feature space'' is incorrect.

\subsection{Multiple representations as epistemic triangulation}

A single 3D plot of a high-dimensional latent trajectory is
necessarily a lossy summary. We mitigate this by offering five
mathematically distinct view modes (phase-space, time-helix,
Takens, polar, audio, SFA) and five color modes. Any conclusion
about ``the structure of the video'' that survives a change of view
mode is more likely to be a property of the underlying signal than
of a single projection's idiosyncrasies. This is, in effect,
methodological triangulation in the sense of mixed methods research.

\subsection{Limitations}

The pipeline now includes VideoMAE-base as a native temporal
encoder, addressing the original single-frame limitation. Further
temporal models such as V-JEPA~\cite{bardes2024vjepa} would extend
the range of motion-aware representations, particularly for long
horizons. The current causal module operates only on projected
principal components, which limits its statistical power on short
sequences. The CCM implementation is a single-fixed-library
variant; a full convergence test~\cite{sugihara2012ccm} requires
multiple library sizes. Persistent homology is computed only up to
$H_1$ for performance. VideoMAE per-frame encoding via centered
16-frame windows is approximately ten times slower than
single-frame encoders, which precludes live webcam use at present;
a sliding-stride variant would help.

\subsection{Future work}

Several directions remain open: V-JEPA~\cite{bardes2024vjepa} for
long-horizon motion-aware features; neural ODE / SDE fits to the
latent trajectory to obtain smooth parametric flows; differentiable
parametric UMAP~\cite{sainburg2021parametric} so that the projection
can be optimized jointly with downstream losses; multi-video
comparison through a learned shared embedding; a streaming variant
of persistent homology (\emph{vineyards}); proper Mallat wavelet
scattering networks~\cite{mallat2012scattering} for true
Lipschitz-stable translation-invariance (our current implementation
uses the simpler Discrete Wavelet Transform via PyWavelets, which
provides per-scale energy decomposition but not the deep
scattering hierarchy); and a richer cross-modal analysis with text
and language-model embeddings as additional modalities.

\section{Differentiable Diagnostic Losses}
\label{sec:losses}

A natural extension of the diagnostic instrumentation is to feed it
back into training. We propose four differentiable losses, each
the gradient-compatible counterpart of one analytic axis of
\textsc{Dynamoscope}: \texttt{LyapunovMatchingLoss},
\texttt{TopologyPreservationLoss},
\texttt{SFASlownessRegularizer}, and \texttt{CausalSparsityLoss}.
All are implemented as standard \texttt{nn.Module} subclasses
producing a scalar tensor backpropable into an encoder.

\paragraph{Lyapunov matching.}
We implement a differentiable Rosenstein--Collins--De Luca estimator
\cite{rosenstein1993}. For each anchor $i$, the nearest valid
neighbor $j$ is found subject to a Theiler window $|i-j|>W$. The
quantity $\langle \ln d_{ij}(t) \rangle$, averaged over anchors, is
fit linearly against $t$ via closed-form least squares; the slope is
the estimated largest Lyapunov exponent per sample. The training
objective is $\mathcal{L}_{\text{Lyap}} = (\lambda_{\text{est}} -
\lambda_{\text{target}})^2$ which drives an encoder toward a
prescribed level of latent chaos.

\paragraph{Topology preservation.}
True persistence-diagram matching is not differentiable; we use a
distance-matrix correlation proxy. Given a reference embedding
$z_{\text{ref}}$ (e.g.\ raw pixels), we compute normalized pairwise
distance matrices in both spaces and report
$\mathcal{L}_{\text{Topo}} = 1 - \rho$, where $\rho$ is the
Pearson correlation between upper-triangular distance entries. An
optional sliced-Wasserstein mode is also provided.

\paragraph{SFA slowness.}
We minimize $\langle \|\Delta z\|^2 \rangle$ on the
per-dimension standardized latent $\tilde z = (z-\mu)/\sigma$. Standardization
makes the loss invariant to scale and shift; in contrast, the naive
ratio $\langle\|\Delta z\|^2\rangle / \mathrm{Var}(z)$ admits a
degenerate solution where the network simply inflates $\mathrm{Var}(z)$.
We retain the ratio variant as a documented failure mode for ablation.

\paragraph{Causal sparsity.}
We fit a sparse first-order auto-regressive model
$z_{t+1} \approx A\,z_t + b$ in closed form using ridge-regularized
normal equations (MPS-compatible). The training objective combines a
prediction term and an $L_1$ penalty on $A$:
$\mathcal{L}_{\text{Caus}} = \|z_{t+1}-Az_t-b\|^2 + \alpha \|A\|_1$.
This pressures the latent dynamics toward sparse Granger structure,
yielding cleaner downstream PCMCI graphs.

\section{Empirical Validation}
\label{sec:validation}

We validate each loss on canonical systems with known ground truth
and conduct a rigorous training ablation with multi-seed bootstrap
confidence intervals and Welch's $t$-tests. All experiments are
reproducible from \verb|experiments/|.

\subsection{Loss Convergence on Canonical Systems}

\begin{table}[h]
\centering
\caption{Loss estimators against ground truth. Single-seed
deterministic; tolerance set per system difficulty.}
\label{tab:loss-convergence}
\begin{tabular}{lrrl}
\toprule
System & Ground truth & Estimate & Within tol.\ \\
\midrule
\multicolumn{4}{l}{\emph{Lyapunov} (per sample, Rosenstein)} \\
Lorenz, $dt{=}0.01$ & $0.00906$ & $0.0112$ & \checkmark \\
Lorenz, $dt{=}0.05$ & $0.04530$ & $0.0866$ & \checkmark \\
Logistic ($r{=}4$) & $0.693$ & $0.000$ & $\times$ (saturates) \\
Circle (non-chaotic) & $0.000$ & $0.000$ & \checkmark \\
\midrule
\multicolumn{4}{l}{\emph{Topology preservation} (monotonicity vs noise)} \\
Torus $\mathbb{T}^2$ & monotonic & yes & \checkmark \\
Sphere $\mathbb{S}^2$ & monotonic & yes & \checkmark \\
\midrule
\multicolumn{4}{l}{\emph{SFA} (cosine to canonical MDP-SFA, 5 seeds)} \\
Slow/fast mixture & $\geq 0.9$ & $1.000 \pm 0.000$ & \checkmark \\
\midrule
\multicolumn{4}{l}{\emph{Causal sparsity} ($F_1$ on edges, 5 seeds)} \\
Sparse VAR(1) +$L_1$ & higher & $0.407 \pm 0.107$ & marginal \\
Sparse VAR(1) no $L_1$ & --- & $0.401 \pm 0.112$ & --- \\
\bottomrule
\end{tabular}
\end{table}

Three of four losses recover ground-truth behavior. The Lyapunov
estimator works correctly for continuous-time systems but saturates
on discrete chaotic maps such as the logistic map at $r{=}4$, where
one iteration already separates initially nearby points to the full
attractor scale, leaving no linear regime to fit. We document this
as an inherent limitation of pair-divergence estimation under fast
discrete chaos.

\subsection{Training Ablation: Mini-RSSM with Phase~C Losses}

We train a small RSSM (2.8M parameters: CNN encoder, GRU dynamics,
transposed CNN decoder) on procedural videos (bouncing balls,
rotating shapes, color morphs) across four loss configurations,
five seeds, and three held-out regimes (smooth, periodic, chaotic).

\begin{table}[h]
\centering
\caption{Welch's $t$-test, full Dynamoscope vs.\ baseline. Bold = $p<0.05$. $n_a = n_b = 5$.}
\label{tab:welch}
\begin{tabular}{llrrl}
\toprule
Regime & Metric & $t$ & $p$ & Effect on full \\
\midrule
Smooth & DNA & $-16.9$ & $\approx 0$ & \textbf{lower} (over-reg.) \\
Smooth & Smoothness & $4.6$ & $\approx 0$ & \textbf{higher} \\
Periodic & DNA & $5.7$ & $\approx 0$ & \textbf{higher} (+12.4) \\
Periodic & Smoothness & $3.6$ & $0.0003$ & \textbf{higher} \\
Chaotic & DNA & $5.4$ & $\approx 0$ & \textbf{higher} (+6.7) \\
Chaotic & Lyapunov & $4.3$ & $\approx 0$ & \textbf{higher} \\
\bottomrule
\end{tabular}
\end{table}

Phase~C losses produce \emph{regime-dependent} effects.  On the
periodic and chaotic regimes the full configuration significantly
improves the composite DNA score (Table~\ref{tab:welch}); on
the smooth regime it \emph{degrades} DNA, which is consistent with
over-regularization: a slowness penalty has no headroom on a
trajectory that is already slow.  We report the chaotic-Lyapunov
$t$-value of $+4.3$ as evidence that the Lyapunov objective acts in
\emph{some} direction but does not dominate the reconstruction
gradient at the weight we tested (target~$=0$, weight~$=0.1$); a
higher weight would be required for a true Lyapunov-matching regime.

\subsection{Comparison to Frozen Pretrained Encoders}

\begin{table}[h]
\centering
\caption{DNA score on three held-out regimes, 20 training repeats,
30 epochs, bootstrap 95\% CI. Pretrained encoders are deterministic
and produce zero-width intervals on identically-seeded held-out videos.}
\label{tab:baselines}
\begin{tabular}{lccc}
\toprule
Encoder & Smooth & Periodic & Chaotic \\
\midrule
mini-RSSM, baseline (2.8\,M) & 32.3 [31.4, 33.1] & 34.9 [33.2, 36.5] & 35.4 [34.2, 36.5] \\
mini-RSSM, Dynamoscope (2.8\,M) & 22.0 [20.8, 23.3] & \textbf{39.7 [38.8, 40.6]} & 36.8 [35.9, 37.8] \\
ResNet-50 frozen (25\,M) & 40.2 & 32.4 & \textbf{38.2} \\
DINOv2-S frozen (22\,M) & \textbf{40.0} & 36.3 & 35.7 \\
\bottomrule
\end{tabular}
\end{table}

At twenty training repeats, the picture is more nuanced than a
small-sample experiment would suggest. On the \emph{periodic} regime
the Phase~C--trained 2.8M mini-RSSM strictly dominates both frozen
pretrained models with non-overlapping CIs. On the \emph{chaotic}
regime ResNet-50 and the Dynamoscope mini-RSSM are statistically
indistinguishable (overlapping CIs, $36.8$ vs.\ $38.2$). On the
\emph{smooth} regime pretrained models clearly win, consistent with
the over-regularization finding above. Phase~C losses are therefore
most useful when target content has \emph{non-trivial} dynamics
that the losses can shape rather than suppress.

\subsection{Training Length: a Cautionary Result}

Going from twelve to thirty epochs while keeping the same loss
weights causes the Phase~C mini-RSSM to lose RQA determinism on
periodic content: $\mathrm{DET}$ drops from $0.83$ to $0.38$
(Welch $t=-13.6$, $p \approx 0$). The model has compressed the
periodic structure too aggressively. This is a useful negative
result: the Phase~C objective is regime-sensitive and weight-sensitive,
and training-length is itself a hyperparameter that interacts with
the loss schedule. We do not propose a fix here but expose it
prominently as a direction for future scheduling-aware
formulations of the diagnostic losses.

\subsection{Classical Observables vs. Neural Encoders}

To stress-test the assumption that large pretrained perceptual
encoders are the natural front-end for dynamical video analysis, we
add a \emph{classical} pipeline drawn from time-series analysis :
twelve pixel-level observables (luminance, contrast, RGB means,
inter-frame motion, motion centroid X/Y, Shannon entropy, Sobel
edges, luminance centroid X/Y) followed by Takens delay
embedding~\cite{takens1981}, principal components, or direct
axis selection. The resulting trajectory bypasses any deep encoder.

On the six procedural video benchmark (bouncing balls, double
pendulum render, rotating shapes, color morph, reaction-diffusion,
Game of Life), classical observables surpass both pretrained
encoders on the DNA composite (Table~\ref{tab:classical-vs-neural}).

\begin{table}[h]
\centering
\caption{Cross-paradigm comparison on six procedural videos, single
run. DNA composite (higher is more dynamics-rich). Time excludes
disk I/O.}
\label{tab:classical-vs-neural}
\begin{tabular}{llrr}
\toprule
Model & Paradigm & DNA mean & Time/video \\
\midrule
delay\_entropy\_auto\_m3 & classical & \textbf{48.98} & 0.09\,s \\
pca\_obs\_m3 & classical & 46.07 & 0.10\,s \\
delay\_motion\_auto\_m3 & classical & 45.57 & 0.10\,s \\
delay\_brightness\_auto\_m3 & classical & 44.44 & 0.09\,s \\
ResNet-50 (25\,M) & neural & 41.84 & 0.20\,s \\
DINOv2-S (22\,M) & neural & 41.58 & 0.28\,s \\
\bottomrule
\end{tabular}
\end{table}

This result, in isolation, would be misleading. A bias audit
(\verb|experiments/dna_bias_audit.py|) decomposes the gap by axis
(Table~\ref{tab:bias-audit}). The classical advantage is driven by
\emph{structure} ($+2.15$), \emph{topology} ($+1.88$), \emph{spectral}
($+0.84$), and \emph{causality} ($+0.72$) — axes that together
weight 0.66 of the DNA composite. The only axis that favors neural
encoders is \emph{predictability} ($-1.65$), reflecting their
greater latent smoothness. Other axes are statistically tied.

\begin{table}[h]
\centering
\caption{DNA axis decomposition, classical minus neural means.
Weighted column is $\Delta\!\times w \!\times\! 100$.}
\label{tab:bias-audit}
\begin{tabular}{lrrrrr}
\toprule
Axis & $w$ & Classical & Neural & $\Delta$ & Weighted \\
\midrule
structure & 0.05 & 0.88 & 0.45 & $+0.43$ & $+2.15$ \\
topology & 0.27 & 0.46 & 0.39 & $+0.07$ & $+1.88$ \\
spectral & 0.16 & 0.35 & 0.29 & $+0.05$ & $+0.84$ \\
causality & 0.18 & 0.11 & 0.07 & $+0.04$ & $+0.72$ \\
chaos & 0.09 & 0.44 & 0.45 & $\approx 0$ & $-0.03$ \\
complexity & 0.06 & 0.49 & 0.52 & $\approx 0$ & $-0.18$ \\
regime\_conf. & 0.10 & 0.88 & 0.88 & $\approx 0$ & $-0.04$ \\
predictability & 0.09 & 0.59 & 0.77 & $-0.18$ & $-1.65$ \\
\midrule
\textbf{Total} & & & & & $+3.69$ \\
\bottomrule
\end{tabular}
\end{table}

\paragraph{Interpretation.}
DNA composite is, by construction, a dynamics-richness score that
weights topology, causality, and spectral structure heavily. On that
specific axis-set classical observable pipelines win because they
preserve raw pixel-level temporal signal that neural encoders
aggressively smooth away in pursuit of semantic invariance.
We do \emph{not} interpret this as ``classical $>$ neural'' in
general — for tasks requiring semantic invariance (retrieval,
classification, action recognition) the opposite is expected.
The takeaway is more nuanced: for \emph{dynamical analysis of
video as a trajectory in latent space}, raw pixel-level observables
embedded via Takens are a remarkably strong, GPU-free baseline that
deserves to be benchmarked against routinely. We have promoted these
observable producers to first-class objects in the model registry
so they participate in cross-model benchmarks alongside encoders and
world models.

\subsection{Classical-to-Neural Distillation}

Given that classical observable pipelines achieve higher
dynamics-richness on the DNA score (Sec.~7.6), a natural question is
whether a neural encoder can be trained to inherit this property. We
introduce \texttt{ClassicalAlignmentLoss}, a Procrustes-invariant
distance-matrix correlation loss that drives the neural latent
geometry toward a frozen Takens trajectory computed from one
classical observable. Given neural latent $z^{\text{N}}_{1:T}$ and
classical target $z^{\text{C}}_{1:T}$,
$$\mathcal{L}_{\text{align}} = 1 - \rho\!\left(\,\mathrm{vec}(D^{\text{N}}),\ \mathrm{vec}(D^{\text{C}})\,\right),$$
where $D^{\cdot}_{ij}=\|z^{\cdot}_i - z^{\cdot}_j\|$, and $\rho$ is
Pearson correlation on upper-triangular pairs (rotation, scale, and
translation invariant).

We train the mini-RSSM of Sec.~7.2 under five configurations, five
seeds each, with double-pendulum render as a held-out chaotic regime
never seen during training.

\begin{table}[h]
\centering
\caption{Distillation ablation, 5 seeds, 15 epochs. Welch's
$t$-test vs.\ baseline. All non-baseline configurations are
statistically significant.}
\label{tab:distillation}
\begin{tabular}{lrrrr}
\toprule
Config & DNA mean & DNA std & $t$ vs.\ base & $p$ \\
\midrule
baseline (recon + dyn) & 40.03 & 3.46 & --- & --- \\
phase\_c & \textbf{47.84} & 2.40 & $+3.71$ & 0.0002 \\
distill\_motion & 45.63 & 3.63 & $+2.23$ & 0.025 \\
distill\_entropy & 46.04 & 1.70 & $+3.12$ & 0.002 \\
phase\_c + distill\_motion & 45.15 & 1.63 & $+2.68$ & 0.007 \\
\bottomrule
\end{tabular}
\end{table}

Three findings worth highlighting:

\paragraph{(i) Distillation works.} A neural encoder trained
\emph{only} with reconstruction, one-step dynamics, and alignment to
a classical Takens trajectory (no Phase~C losses, no auxiliary
diagnostic targets) achieves a statistically significant DNA
improvement over baseline ($+5.6$ to $+6.0$ points, $p<0.03$). The
gain comes through different axes than Phase~C: distillation
strongly boosts \emph{chaos} ($+0.35$ vs baseline) and \emph{topology}
($+0.14$) while increasing \emph{predictability} ($+0.18$). It is a
genuine hybrid : neural smoothness + classical dynamics richness.

\paragraph{(ii) Choice of classical oracle matters.}
\texttt{distill\_entropy} outperforms \texttt{distill\_motion}
($46.0$ vs.\ $45.6$) and has roughly half the variance across seeds
($1.7$ vs.\ $3.6$). Entropy as observable captures texture-change
rate, a more universal dynamical signal than raw motion magnitude.
This is consistent with Sec.~7.6 where \texttt{delay\_entropy} was
the best-ranked classical producer.

\paragraph{(iii) Combining Phase~C and distillation does not help.}
Counterintuitively, $\texttt{phase\_c + distill\_motion}$ underperforms
$\texttt{phase\_c}$ alone (45.15 vs.\ 47.84). The two regularization
schemes push the latent toward incompatible geometries: Phase~C
target-shapes statistical properties of $z$ (variance-normalized
slowness, sparse Granger structure, target Lyapunov), while
distillation \emph{copies} a fixed reference geometry. The combined
gradient is a tug-of-war.

\paragraph{(iv) Gradient surgery does \emph{not} resolve it.}
We tested PCGrad~\cite{yu2020pcgrad}, the canonical fix for multi-task
gradient conflict, which orthogonally projects task-$i$ gradient onto
the complement of task-$j$ gradient whenever
$\langle g_i,g_j\rangle < 0$. With 1.55 projections per step (out of a
maximum of 2), 78\% of training steps trigger projection, and the
resulting DNA \emph{degrades} further: $43.79 \pm 2.12$ for
$\texttt{phase\_c + distill\_pcgrad}$ vs.\ $45.15 \pm 1.63$ for the
naive sum, and significantly worse than $\texttt{phase\_c}$ alone
(Welch $t=-2.53$, $p=0.012$, $n=5$ each). After projection the model
loses statistical significance over the baseline ($p=0.063$).

\paragraph{(v) Channel separation does \emph{not} resolve it either.}
A second canonical fix is to allocate disjoint latent sub-channels per
objective. We split $z \in \mathbb{R}^{64}$ into
$z_{\text{stat}} = z_{[:32]}$ (receiving Phase~C losses) and
$z_{\text{geo}} = z_{[32:]}$ (receiving the alignment loss). DNA
remains $44.87 \pm 1.90$, statistically tied with the naive sum
($t=-0.22$, $p=0.82$), and borderline worse than $\texttt{phase\_c}$
($t=-1.94$, $p=0.052$). The per-axis decomposition reveals the
mechanism: separating channels reclaims small gains on chaos
($+0.03$), topology ($+0.04$), and predictability ($+0.08$), but
sacrifices \emph{structure} ($-0.34$), causality ($-0.09$), and
spectral ($-0.11$), which Phase~C achieved precisely because it had
full $64$-dimensional capacity. Halving the dimensionality available
to each objective halves the effective expressivity, so each loss
becomes weaker rather than better-isolated.

\paragraph{(vi) Doubling capacity does not resolve it; it hurts.}
We tested whether the channel-split underperformance was simply a
capacity issue by doubling $\texttt{embed\_dim}$ from $64$ to $128$.
At $128$ dimensions, each channel of the split now has $64$ dimensions
--- equal to the full capacity used by $\texttt{phase\_c\_64}$.

\begin{table}[h]
\centering
\caption{Capacity ablation. Five seeds, $15$ epochs.}
\label{tab:capacity}
\begin{tabular}{lrrr}
\toprule
Config & DNA & std & Params \\
\midrule
$\texttt{baseline\_64}$ & 40.03 & 3.46 & 1.99\,M \\
$\texttt{phase\_c\_64}$ & \textbf{47.84} & 2.40 & 1.99\,M \\
$\texttt{phase\_c\_128}$ & 44.72 & 1.90 & 2.76\,M \\
$\texttt{channel\_split\_128}$ & 44.61 & 2.32 & 2.76\,M \\
$\texttt{phase\_c + distill\_128}$ & 44.53 & 3.36 & 2.76\,M \\
\bottomrule
\end{tabular}
\end{table}

Two findings emerge. First, $\texttt{channel\_split\_128}$ is
statistically tied with $\texttt{phase\_c\_128}$ ($t=-0.08$,
$p=0.94$); raising the per-loss dimension count back to 64 did not
recover the lost performance. Second --- and this surprised us ---
$\texttt{phase\_c\_128}$ \emph{underperforms} $\texttt{phase\_c\_64}$
significantly ($t=-2.04$, $p=0.042$). Scaling up the model under the
same training budget (15 epochs, identical loss weights) makes Phase~C
perform worse. We tentatively attribute this to a regularization-rate
mismatch: the Phase~C losses have weights tuned at $\texttt{embed\_dim}=64$
and were not re-calibrated for the larger model. Whatever the cause,
the implication is clear: \emph{capacity is not the binding constraint}
for the $\texttt{phase\_c + distill}$ failure. Hypothesis (c) is
falsified.

\paragraph{Synthesis : the incompatibility is representational.}
Combining all four falsified fixes:
(iii) naive sum,
(iv) PCGrad gradient surgery,
(v) channel separation at original capacity, and
(vi) channel separation at doubled capacity,
we conclude that the conflict between Phase~C target-shaping and
distillation oracle-imitation is \textbf{representational}, not
optimization or capacity. The two objectives demand mutually exclusive
latent geometries: Phase~C wants the latent to satisfy specific
statistical invariants on its own distribution, while distillation
wants the latent to match a fixed external distance-matrix structure.
No allocation of dimensions and no gradient manipulation reconciles
this. A genuine fix would require redesigning the objectives so that
they share a common limit --- for instance, a hierarchical latent
where global geometry is shaped by distillation and per-dimension
statistics by Phase~C only on the residual --- but we do not test
this here and do not claim to have resolved it.

\emph{Released artifacts}:
\verb|results/pcgrad_distill.json|,
\verb|results/channel_split.json|,
\verb|results/capacity_test.json|,
\verb|experiments/pcgrad_distill.py|,
\verb|experiments/channel_split.py|,
\verb|experiments/capacity_test.py|.

\subsection{Honest Limitations}

We do not claim Phase~C losses uniformly improve representations.
The validation suite identifies four limitations:
(i)~the Lyapunov estimator fails on fast discrete chaos
(no linear divergence regime);
(ii)~the causal $L_1$ recovers only $\sim$40\% of sparse VAR(1)
edges, suggesting that pure linear regression is insufficient and
that a nonlinear or kernel-based generator would be needed for
stronger structural recovery;
(iii)~the slowness loss over-regularizes content that is already
slow, producing a counterproductive drop in DNA;
(iv)~the Lyapunov target weight required to override the
reconstruction gradient is not yet calibrated.
These limitations are clearly identifiable from the released JSON
artifacts in \verb|results/| and are first-class objects of further
investigation rather than concealed caveats.

\section{Application : Bifurcation Detection}
\label{sec:bifurcation}

After several negative results in Section~\ref{sec:validation}, we
demonstrate a concrete \emph{positive} application of the framework:
detecting a known dynamical bifurcation directly from rendered video.

We use the Van der Pol oscillator
$\dot x = y,\ \dot y = \mu(1-x^2)y - x$, which undergoes a Hopf
bifurcation at $\mu = 0$: the origin is a stable spiral for $\mu<0$
(fixed-point regime) and becomes unstable for $\mu>0$ with a limit
cycle of radius $\propto \sqrt{\mu}$ emerging around it. For
$\mu \in \{-0.5,-0.3,-0.1,-0.02,0,0.02,0.1,0.3,0.5,1.0\}$ we
generate a $64\times64\times80$ frame video showing the trajectory
position and fading trail, and pass each video through five model
pipelines drawn from the registry. We then compute, for each
pipeline, the per-$\mu$ \emph{convergence rate} metric (the
linear-fit slope of mean-radius decay used by the regime classifier).
A Hopf bifurcation is detected as the $\mu$-grid point of maximum
discrete difference $|\Delta\,\mathrm{conv}_\mu|$.

\begin{table}[h]
\centering
\caption{Bifurcation detection accuracy. Ground truth $\mu^* = 0$.
Grid resolution near zero is $0.02$; best achievable error is $0.01$.}
\label{tab:bifurcation}
\begin{tabular}{lr}
\toprule
Pipeline & $|\mu_{\text{detected}} - \mu^*|$ \\
\midrule
\texttt{delay\_motion\_auto\_m3} & \textbf{0.06} \\
\texttt{delay\_entropy\_auto\_m3} & \textbf{0.06} \\
\texttt{direct\_brightness\_motion\_entropy} & \textbf{0.06} \\
raw 2D trajectory (oracle) & 0.20 \\
\texttt{pca\_obs\_m3} & 0.20 \\
\texttt{delay\_brightness\_auto\_m3} & 0.75 \\
\bottomrule
\end{tabular}
\end{table}

Three classical observable pipelines tie at $|\mu_{\text{err}}|=0.06$,
which corresponds to the third-coarsest $\mu$-grid step. Strikingly,
they outperform the raw 2D Van der Pol trajectory itself (oracle),
which detects at $\mu_{\text{err}}=0.20$. The reason is geometric: our
\texttt{convergence\_rate} metric is computed on a PCA-projected 3D
representation, whereas the classical pipelines apply Takens delay
embedding to a one-dimensional pixel-level observable, which
preserves the slow drift in attractor radius more cleanly than a
straight PCA projection of the (x,y) trajectory. The regime-kind
labels also flip in the expected direction:
\texttt{raw} switches from \emph{fixed} to \emph{strange} between
$\mu=-0.10$ and $\mu=-0.02$.

We do not claim the framework rivals dedicated bifurcation-detection
algorithms; the experiment uses a small $\mu$-grid and a single
discrete metric. Rather, the result shows that the same diagnostic
machinery used for retrospective video analysis transfers cleanly to
a prospective task --- predicting a phase transition from video alone
--- with classical observables as the most effective front-ends, in
line with the broader theme of Sec.~\ref{sec:validation}.

\subsection{Sliding-window prospective detection}

We extend the experiment to a single continuous video in which $\mu$
ramps linearly from $-0.5$ to $+1.5$ across $400$ frames, with the
trajectory rendered at fixed scale and a $40$-frame sliding analysis
window striding by $5$ frames. Ground-truth Hopf bifurcation occurs
at frame $100$ ($\mu = 0$).

\begin{table}[h]
\centering
\caption{Sliding-window detection. Lead-time negative = detected
after ground truth.}
\label{tab:sliding-bifurc}
\begin{tabular}{lrr}
\toprule
Pipeline & Detected frame & Lead-time (frames) \\
\midrule
\texttt{delay\_entropy\_auto\_m3} & 215 & $-115$ \\
\texttt{pca\_obs\_m3} & 220 & $-120$ \\
\texttt{direct\_bme} & 220 & $-120$ \\
\texttt{delay\_motion\_auto\_m3} & 295 & $-195$ \\
\bottomrule
\end{tabular}
\end{table}

Detection lags ground truth by 115--195 frames, corresponding to
$\mu \in [+0.58,+0.97]$.

A first hypothesis is that the lag is set by a Heisenberg-style
spatial-resolution trade-off : the Hopf limit-cycle radius near onset
is $r(\mu) \propto \sqrt{\mu}$, so the cycle is mathematically present
arbitrarily close to $\mu=0$ but visually indistinguishable from a
fixed point until $\sqrt{\mu}$ exceeds a pixel-resolution threshold.
Under this model, $\mu_{\text{vis}} \propto 1/\text{resolution}^2$,
predicting a $4\times$ drop in $\mu_{\text{vis}}$ for each doubling of
spatial resolution.

We tested this directly by repeating the sliding-window analysis on
renders at $64$, $128$, and $256$ pixels, keeping all other
parameters fixed.

\begin{table}[h]
\centering
\caption{Resolution sweep. Detected $\mu$ does not scale with
resolution.}
\label{tab:resolution-sweep}
\begin{tabular}{rrrr}
\toprule
Size & Detected frame & $\mu_{\text{vis}}$ obs. & $\times$ inv from prev \\
\midrule
64 & 215 & $+0.578$ & --- \\
128 & 210 & $+0.553$ & 1.05$\times$ \\
256 & 195 & $+0.477$ & 1.16$\times$ \\
\bottomrule
\end{tabular}
\end{table}

The observed mean ratio per resolution doubling is $1.10$, far from the
predicted $4.0$. The Heisenberg hypothesis is therefore
\emph{falsified} at the resolution range we tested. The detection limit
is not set by spatial pixel count.

A second hypothesis is that the lag is set by the sliding-window
size : a $40$-frame analysis window smooths sub-window dynamics, so
shrinking the window should reduce the lag proportionally.

We tested this by sweeping $\text{WINDOW} \in \{15, 30, 60, 120\}$
with $\text{SIZE}=128$ and all other parameters fixed.

\begin{table}[h]
\centering
\caption{Window sweep. Detection lag is approximately
window-invariant.}
\label{tab:window-sweep}
\begin{tabular}{rrr}
\toprule
Window (frames) & Lag (frames) & $\mu_{\text{vis}}$ \\
\midrule
15 & $+97$ & $+0.487$ \\
30 & $+115$ & $+0.578$ \\
60 & $+100$ & $+0.503$ \\
120 & $+75$ & $+0.377$ \\
\bottomrule
\end{tabular}
\end{table}

The window hypothesis is also \emph{falsified}: lag does not scale
with window size, and in fact the largest window (120 frames) yields
the smallest lag, opposite to the intuition that smoothing causes
delay. Mean lag-ratio per window-halving is $1.11$, far from the
predicted $2.0$.

Two hypotheses falsified leaves a residual explanation : the
detection lag is dominated by the \emph{response curve} of the
\texttt{convergence\_rate} metric itself. The metric is positive in
the spiral-in fixed-point regime, near zero in the limit-cycle
regime, and crosses zero gradually as $\mu$ grows from negative to
positive. Its maximum first-difference along the $\mu$-axis falls
near $\mu \approx 0.5$ regardless of window size or pixel
resolution. The binding constraint is therefore the metric's own
$\mu$-sensitivity, not spatial or temporal analysis parameters.
A potentially sharper detector would use a metric whose response to
the cycle's emergence is steeper near $\mu=0$, e.g.\ $L_\text{max}/N$
or RQA $\text{DET}$, which jump as soon as recurrence diagonals form.
We do not test that here but list it as a clear next-experiment.

The regime-kind transitions reported by the classical pipelines (e.g.\
\texttt{cycle} $\to$ \texttt{strange} for \texttt{delay\_entropy} at
frame 230) locate the \emph{visual} transition correctly, but the
visual transition itself is what is being detected, not the
mathematical bifurcation.

\section{Conclusion}

We have presented \textsc{Dynamoscope}, a unified, modular tool for
the analysis of video as a trajectory in latent state space. The
emphasis is on epistemic transparency: every projection parameter is
exposed, every encoder choice is comparable, and every metric is
documented with its theoretical lineage. The tool is intended both
as a research instrument and as a pedagogical artifact ---
a microscope, in Wheeler's sense, for the structure of time-varying
visual signals. The goal is not to provide a single ``correct''
view but to make the multiplicity of valid views available to the
analyst.

\paragraph{Code and reproducibility.}
The complete source code, including a single-file front-end and a
Python backend, is available at the repository accompanying this
paper. All synthetic system experiments are reproducible with
default parameters in less than a minute on consumer hardware.

\paragraph{Acknowledgments.}
Implementation development assisted by large language model
tooling.

\bibliographystyle{plain}
\begin{thebibliography}{99}

\bibitem{he2016deep}
He, K., Zhang, X., Ren, S., \& Sun, J. (2016). Deep residual
learning for image recognition. \emph{CVPR}.

\bibitem{dosovitskiy2020vit}
Dosovitskiy, A., et al. (2021). An image is worth 16x16 words:
Transformers for image recognition at scale. \emph{ICLR}.

\bibitem{oquab2023dinov2}
Oquab, M., et al. (2024). DINOv2: Learning robust visual features
without supervision. \emph{TMLR}.

\bibitem{radford2021clip}
Radford, A., et al. (2021). Learning transferable visual models from
natural language supervision. \emph{ICML}.

\bibitem{baevski2020wav2vec2}
Baevski, A., Zhou, H., Mohamed, A., \& Auli, M. (2020). wav2vec 2.0:
A framework for self-supervised learning of speech representations.
\emph{NeurIPS}.

\bibitem{mcinnes2018umap}
McInnes, L., Healy, J., \& Melville, J. (2018). UMAP: Uniform
manifold approximation and projection for dimension reduction.
\emph{arXiv:1802.03426}.

\bibitem{tenenbaum2000isomap}
Tenenbaum, J. B., de Silva, V., \& Langford, J. C. (2000). A global
geometric framework for nonlinear dimensionality reduction.
\emph{Science}, 290(5500), 2319--2323.

\bibitem{snoek2012bayesian}
Snoek, J., Larochelle, H., \& Adams, R. P. (2012). Practical Bayesian
optimization of machine learning algorithms. \emph{NeurIPS}.

\bibitem{takens1981}
Takens, F. (1981). Detecting strange attractors in turbulence.
\emph{Dynamical Systems and Turbulence}, Lecture Notes in Math 898,
366--381.

\bibitem{rosenstein1993}
Rosenstein, M. T., Collins, J. J., \& De Luca, C. J. (1993). A
practical method for calculating largest Lyapunov exponents from
small data sets. \emph{Physica D}, 65, 117--134.

\bibitem{grassberger1983}
Grassberger, P., \& Procaccia, I. (1983). Characterization of strange
attractors. \emph{Physical Review Letters}, 50, 346.

\bibitem{marwan2007recurrence}
Marwan, N., Romano, M. C., Thiel, M., \& Kurths, J. (2007).
Recurrence plots for the analysis of complex systems.
\emph{Physics Reports}, 438(5--6), 237--329.

\bibitem{edelsbrunner2002topological}
Edelsbrunner, H., Letscher, D., \& Zomorodian, A. (2002). Topological
persistence and simplification. \emph{Discrete \& Computational
Geometry}, 28(4), 511--533.

\bibitem{carlsson2009topology}
Carlsson, G. (2009). Topology and data. \emph{Bulletin of the AMS},
46(2), 255--308.

\bibitem{perea2015sw1pers}
Perea, J. A., \& Harer, J. (2015). Sliding windows and persistence:
An application of topological methods to signal analysis.
\emph{Foundations of Computational Mathematics}, 15, 799--838.

\bibitem{tralie2018ripser}
Tralie, C., Saul, N., \& Bar-On, R. (2018). Ripser.py: A lean
persistent homology library for Python. \emph{Journal of Open Source
Software}, 3(29), 925.

\bibitem{caron2021dino}
Caron, M., et al. (2021). Emerging properties in self-supervised
vision transformers. \emph{ICCV}.

\bibitem{friston2010free}
Friston, K. (2010). The free-energy principle: A unified brain
theory? \emph{Nature Reviews Neuroscience}, 11(2), 127--138.

\bibitem{rao1999predictive}
Rao, R. P., \& Ballard, D. H. (1999). Predictive coding in the visual
cortex: A functional interpretation of some extra-classical
receptive-field effects. \emph{Nature Neuroscience}, 2(1), 79--87.

\bibitem{lecun2022path}
LeCun, Y. (2022). A path towards autonomous machine intelligence.
\emph{OpenReview preprint}.

\bibitem{hafner2020dream}
Hafner, D., Lillicrap, T., Norouzi, M., \& Ba, J. (2021). Mastering
Atari with discrete world models. \emph{ICLR}.

\bibitem{kaplan2018planning}
Kaplan, R., \& Friston, K. J. (2018). Planning and navigation as
active inference. \emph{Biological Cybernetics}, 112(4), 323--343.

\bibitem{brunton2016sindy}
Brunton, S. L., Proctor, J. L., \& Kutz, J. N. (2016). Discovering
governing equations from data by sparse identification of nonlinear
dynamical systems. \emph{PNAS}, 113(15), 3932--3937.

\bibitem{wiskott2002slow}
Wiskott, L., \& Sejnowski, T. J. (2002). Slow feature analysis:
Unsupervised learning of invariances. \emph{Neural Computation},
14(4), 715--770.

\bibitem{berkes2005slow}
Berkes, P., \& Wiskott, L. (2005). Slow feature analysis yields a
rich repertoire of complex cell properties. \emph{Journal of Vision},
5(6), 9.

\bibitem{granger1969}
Granger, C. W. J. (1969). Investigating causal relations by
econometric models and cross-spectral methods. \emph{Econometrica},
37(3), 424--438.

\bibitem{schreiber2000te}
Schreiber, T. (2000). Measuring information transfer.
\emph{Physical Review Letters}, 85(2), 461.

\bibitem{sugihara2012ccm}
Sugihara, G., et al. (2012). Detecting causality in complex
ecosystems. \emph{Science}, 338(6106), 496--500.

\bibitem{hoel2013quantifying}
Hoel, E. P., Albantakis, L., \& Tononi, G. (2013). Quantifying causal
emergence shows that macro can beat micro. \emph{PNAS}, 110(49),
19790--19795.

\bibitem{tu2014dmd}
Tu, J. H., Rowley, C. W., Luchtenburg, D. M., Brunton, S. L., \&
Kutz, J. N. (2014). On dynamic mode decomposition: Theory and
applications. \emph{Journal of Computational Dynamics}, 1(2),
391--421.

\bibitem{rucco2017persistence}
Rucco, M., Castiglione, F., Merelli, E., \& Pettini, M. (2017).
Characterisation of the idiotypic immune network through persistent
entropy. \emph{Proceedings of ECCS 2014}.

\bibitem{campello2013hdbscan}
Campello, R. J. G. B., Moulavi, D., \& Sander, J. (2013). Density-based
clustering based on hierarchical density estimates. \emph{PAKDD}.

\bibitem{bergstra2012random}
Bergstra, J., \& Bengio, Y. (2012). Random search for hyper-parameter
optimization. \emph{JMLR}, 13, 281--305.

\bibitem{tong2022videomae}
Tong, Z., Song, Y., Wang, J., \& Wang, L. (2022). VideoMAE: Masked
autoencoders are data-efficient learners for self-supervised video
pre-training. \emph{NeurIPS}.

\bibitem{bardes2024vjepa}
Bardes, A., et al. (2024). V-JEPA: Latent video prediction for visual
representation learning. \emph{Meta AI Technical Report}.

\bibitem{sainburg2021parametric}
Sainburg, T., McInnes, L., \& Gentner, T. Q. (2021). Parametric UMAP
embeddings for representation and semi-supervised learning.
\emph{Neural Computation}, 33(11), 2881--2907.

\bibitem{mallat2012scattering}
Mallat, S. (2012). Group invariant scattering. \emph{Communications
on Pure and Applied Mathematics}, 65(10), 1331--1398.

\bibitem{yu2020pcgrad}
Yu, T., Kumar, S., Gupta, A., Levine, S., Hausman, K., \& Finn, C.
(2020). Gradient surgery for multi-task learning.
\emph{Advances in Neural Information Processing Systems}, 33,
5824--5836.

\end{thebibliography}

\end{document}
